\documentclass[11pt]{article}
\usepackage[margin=1in]{geometry}
\usepackage[T1]{fontenc}
\usepackage{microtype}
\usepackage{amsmath,amssymb}
\usepackage{booktabs,tabularx,array,longtable}
\usepackage{graphicx}
\usepackage[hidelinks]{hyperref}
\usepackage{xurl}
\usepackage{listings}
\setlength{\parindent}{0pt}
\setlength{\parskip}{0.5\baselineskip}
\setlength{\emergencystretch}{2em}
\lstset{basicstyle=\ttfamily\small,breaklines=true,columns=fullflexible,frame=single}
\newcommand{\code}[1]{\expandafter\nolinkurl\expandafter{\detokenize{#1}}}

\title{Is AI ``Slop'' Taking Over the Internet?\\
\large From an Interpretable Pilot to V11: Authorship Detection,\\
Dataset Shift, and the Limits of Benchmark Evidence}
\author{Research manuscript}
\date{September 10, 2026}

\begin{document}
\maketitle

\begin{abstract}
The growth of machine-written text raises a measurement problem before it
raises a classification problem: authorship, quality, and population prevalence
are different targets. This study reviews the evidence motivating that
distinction and reconstructs a sequence of detector-development experiments,
from a handcrafted logistic baseline to V11, a linear support vector classifier
using normalized word and character TF--IDF features. The initial 24-document
historical pilot achieved test AUROC 1.0 but detected none of three generated
articles at its validation-selected threshold. Later experiments on a
300-document modern benchmark improved comparison-set performance, yet the
project reported failures on short external examples. A short-text specialist
increased the comparison-set human false-positive rate to 30\%; formatting
normalization alone did not resolve reported transfer failures. V11 combines
paired HC3 answers with benchmark material, using 2,248 training and 620
validation texts. Its recorded validation accuracy is 0.9742, recall 0.9935,
false-positive rate 0.0452, and AUROC 0.9991. These are selection-set results,
not independent estimates of web-wide performance. A repository audit confirms
corpus construction and metric arithmetic but finds an identical human passage
in both V11 splits, missing external sample files, and incomplete subgroup
result retention. The contribution is an exploratory account of operating-point
failure, dataset dependence, and reproducibility limits. Independently
adjudicated quality detection, untouched external validation, and population
prevalence estimation remain uncompleted.
\end{abstract}

\section{Introduction and Research Questions}
Cheap automated publishing could make useful writing harder to find. That
possibility motivates this study; it is not a conclusion built into its design.
An AI-assisted explanation may be accurate and useful, while a human-written
article may be misleading. A classifier that distinguishes writing styles
does not establish which document deserves trust. Nor does a large volume of
generated text establish that most of the Internet is low-quality output.

The original project proposed a joint authorship-and-quality detector and a
separate prevalence study. Its implemented research has developed differently:
the completed model experiments concern \emph{binary authorship labels}, while
independent quality assessment and probability-based web sampling remain
prospective. The paper preserves that larger motivation but evaluates the
research that exists, rather than presenting the proposal as completed work.
Its central empirical question is whether favorable classification results
survive changes in the source, length, formatting, and composition of text.

Four questions retain the original intellectual direction:
\begin{enumerate}
\item \textbf{RQ1---Prevalence:} Has the proportion of AI-generated or
AI-assisted text increased within a clearly defined web sampling frame?
\item \textbf{RQ2---Quality:} What proportion of that content meets an
independently defined low-quality rubric?
\item \textbf{RQ3---Detection:} Can a transparent model identify the
intersection of machine authorship and low quality while limiting false
accusations?
\item \textbf{RQ4---Generalization:} How does detection change for unfamiliar
publishers, generators, genres, and edited text?
\end{enumerate}
The present experiments address the authorship component of RQ3 and provide
exploratory evidence relevant to RQ4. They do not answer RQ1 or RQ2 with new
population data. This boundary matters: a useful negative result about detector
transfer is not a failed prevalence estimate, and an accurate benchmark
classifier is not automatically a validated ``slop'' detector.

We reconstruct the development sequence from source code, raw and processed
datasets, saved results, model metadata, and documentation. V11 is the current
primary model, not necessarily the best model for an independently specified
deployment population. Earlier failures remain part of the evidence. The
paper distinguishes archived experimental results from computations reproduced
during this revision and from observations retained only in project notes.
The accompanying audit records those distinctions explicitly.

\section{Existing Evidence: What Is Established?}
\subsection{Growth depends on the denominator}
Pew Research Center's August 20, 2026 analysis reports signs of AI authorship
in about 10\% of its July 2026 English-language webpage sample, and in more
than one-third of the subset dated after ChatGPT's release \cite{pew}.
Its sampling covered 10,000 pages from each of 49 Common Crawl snapshots
between January 2021 and July 2026. Publication dates were available for only
about 10--15\% of pages in a crawl sample; dated content is not a random subset
of the web \cite{pewmethods}. The estimates concern detector-indicated
authorship or substantial editing, not independently verified low quality.

Dolezal et al. report that roughly 35\% of newly published websites in their
Internet Archive-based study were classified as AI-generated or AI-assisted
by mid-2025. Their preprint does not find statistically significant support
for a decline in factual accuracy or stylistic diversity associated with
increasing AI text \cite{dolezal}. Those findings do not establish that factual
accuracy is unaffected in every setting; they also do not support treating
machine authorship as a quality label.

These estimates have different denominators and should not be pooled.
Existing pages, newly published sites, and the pages people actually read
are distinct populations. A publisher could produce many little-read pages,
while a single widely viewed page could dominate audience exposure. A claim
about a majority of page counts would therefore still require a separate
argument about exposure. None of the experiments reported below samples
either population probabilistically.

\subsection{Detection is not ground truth}
Published evaluations make transfer a substantive part of detector assessment.
RAID evaluates robustness across generators, domains, decoding choices, and
adversarial changes \cite{raid}. Tufts et al. emphasize detection at specified
false-positive rates and document failures on unfamiliar settings
\cite{tufts}. M4 likewise shows that performance on known domains and generators
does not ensure performance on unseen ones \cite{m4}. These are findings from
external studies, not measurements made with V11.

False positives also need to be examined across writing populations. Liang
et al. report bias against non-native English writing in the detectors they
studied \cite{liang}. The present repository does not contain sufficiently
documented proficiency groups to repeat that assessment. We neither infer
language background from prose nor claim that normalization resolves such
disparities. A detector prediction remains a model output, not evidence that
the author used a particular tool.

Mixed-authorship resources such as CoAuthor and human evaluation resources
such as SummEval illustrate related but distinct research tasks
\cite{coauthor,summeval}. They were considered in the earlier planning
materials; neither supplies the completed quality labels used by this
project, because no such completed quality-label corpus is present. Origin
guesses, writing proficiency ratings, and summary-quality judgments cannot
silently substitute for the rubric defined here.

\section{Operational Definition of AI ``Slop''}
Let $A_i=1$ denote documented fully machine-generated origin and $Q_i=1$
denote independently adjudicated low quality. The proposed joint target is
\begin{equation}
S_i=A_iQ_i.
\label{eq:slop}
\end{equation}
This operational definition gives an informal expression a limited meaning
within the proposed study. It does not identify a universally agreed category,
establish intent, or demonstrate mass production. Mixed human--AI authorship
is recorded separately. Substantially assisted writing could enter a secondary
analysis only under an explicit alternative origin definition.

The retained rubric scores factual and evidential adequacy, information value,
and relevance/coherence on a scale from 0 (no material defect), through 1
(limited defect), to 2 (substantial defect). The low-quality decision is positive
when at least two dimensions receive 2, or when a verified central factual
error defeats the document's purpose. Unsupported evidence and unresolved
evidence are not identical: inadequate information may require an uncertain
label rather than a finding of falsehood. Genre-specific expectations also
matter. Fiction, opinion, and personal narrative cannot be judged by the same
evidential conventions as a factual explanation.

Two independent human raters would assess texts without seeing their origin
labels or detector scores; a third would adjudicate disagreements. Reasons
and evidence references should make decisions reviewable. Language proficiency,
stylistic preference, repeated technical terms, and the absence of academic
citations are not automatic quality defects. The four origin-by-quality cells
must all be represented to prevent the shortcut ``generated means bad.''

This annotation procedure is \emph{not completed}. The benchmark's raw JSON
does contain annotator guesses, confidence judgments, comments, and outputs
from other detectors. They concern perceived authorship, not independently
adjudicated quality under this rubric. The processed quality-label fields
remain empty. V11 therefore learns the datasets' supplied binary origin
labels, not $Q_i$ or $S_i$; their provenance is discussed below rather than
assumed from the presence of a label column.

\section{Study Design: Testing the Takeover Hypothesis}
\subsection{Two distinct datasets}
The original design separates a \emph{prevalence corpus}, sampled from a
defined web population, from a \emph{detector corpus}, constructed to support
supervised learning and error assessment. That distinction still applies.
The historical proposal targeted 4,000 sampled pages across four annual
windows in 2022--2025 and 2,400 independently labeled origin-by-quality
documents. These were collection targets, not achieved samples or a power
analysis. Subsequent development corpora do not fulfill those targets merely
because they contain thousands of answers.

\textbf{The 1880 pilot.} The first experiment pairs 12 excerpts from
\emph{Scientific American}, December 18, 1880, with 12 articles authored by
the assistant preparing the project. The source is identified in the pilot
protocol as Project Gutenberg ebook 21081. The historical issue supports
pre-language-model provenance, subject to transcription caveats. It does not
make the prose representative of contemporary writing or establish the truth
of its claims. The generating assistant had the source passages in context.
The recorded seed controls style assignment, not independent language-model
sampling; the backend identifier and generation settings were not exposed.

\textbf{The 1900 expansion.} A separate source, \emph{The Popular Science
Monthly}, June 1900, is retained with the Gutenberg notice for ebook 47227.
Eight extracted human articles supplied a 173-row human-only chunk corpus.
V3 instead uses 30 selected human excerpts and 30 fixed generated articles,
split into 30 training, 14 validation, and 16 comparison rows. The scripts
\code{make_ai_1900.py} and \code{add_ai_1900_more.py} write literal article
strings; rerunning them is not a new generation experiment. This expansion
is distinct from the original 24-document 1880 pilot.

\textbf{The modern benchmark.} The uploaded
\code{data/benchmark/human_detectors.json} contains 300 records, with 150
human-labeled and 150 AI-labeled texts. It records source, title, issue/date,
link, generation condition, and additional human and detector judgments.
The five conditions are \code{gpt-4o}, \code{paraphrased_gpt-4o},
\code{claude}, \code{o1-pro}, and \code{humanized_o1-pro}. A condition value
on a human row identifies its comparison condition, not a machine generator
for that human text.

The raw records contain 150 distinct source--title--issue article tuples,
each represented once per origin label. They are organized into 30 numeric
\code{prompt_id} groups, each containing five human and five generated rows
across conditions. Thus 30 prompt groups do \emph{not} mean 30 underlying
articles. The grouped CSV assigns the lexicographically sorted prompt IDs
to 18 training, six validation, and six comparison groups, yielding
180/60/60 rows. The retained sources include news organizations and science
and general-interest magazines. They are not a representative mix of all
web genres. Mean document length is 738.15 whitespace-delimited words for
human-labeled text and 724.03 for generated text, computed from the stored CSV.

The benchmark's links support source tracing, but the upload lacks a complete
collection protocol, generation log, upstream version, and item-level reuse
ledger. Its date field mixes recognizable dates with numeric encodings that
have not been consistently resolved. We use the supplied labels for describing
the experiment without claiming an independent provenance adjudication.
Historical filenames beginning \code{hc3_} are misleading here: benchmark
experiments through V8 used this JSON-derived corpus, not the subsequently
acquired HC3 question-and-answer files.

\textbf{HC3 and the combined V11 corpus.} The three raw HC3 files contain
3,933 finance, 1,248 medicine, and 1,187 open-QA question records. An early
extraction retains every answer with at least 80 whitespace-delimited words,
giving 11,593 answers (3,346 human-labeled and 8,247 ChatGPT-labeled). That
imbalanced extraction is not V11's training set.

The paired builder instead retains a question only when both answer lists
contain an eligible answer. It selects the first qualifying human answer and
the first qualifying ChatGPT answer, rather than sampling every answer. The
group key combines domain and raw-file line number. The resulting 3,278
question pairs contain 6,556 texts, exactly balanced by origin. A pandas shuffle
with seed 42 precedes a global 80/20 pair split: 2,622 pairs for training and
656 for validation. The split is not domain-stratified.

V11 caps the number of retained pairs per domain at 600 for training and 150
for validation, sampling complete pairs with seed 42. It then adds the
benchmark's existing training and validation rows. Table~\ref{tab:corpus}
shows the actual composition. Both the raw-to-paired reconstruction and the
paired-to-combined membership reconstruction match the uploaded CSVs in the
revision audit. This verifies the transformation, not provenance, copyright
clearance, or representativeness.

\begin{table}[htbp]
\centering\small
\begin{tabular}{lrrrr}
\toprule
& \multicolumn{2}{c}{Training} & \multicolumn{2}{c}{Validation} \\
Component & Human & AI & Human & AI \\
\midrule
HC3 finance & 600 & 600 & 150 & 150 \\
HC3 medicine & 413 & 413 & 126 & 126 \\
HC3 open QA & 21 & 21 & 4 & 4 \\
Modern benchmark & 90 & 90 & 30 & 30 \\
\midrule
Total & 1,124 & 1,124 & 310 & 310 \\
\bottomrule
\end{tabular}
\caption{Audited V11 composition: 2,248 training and 620 validation texts.
Only finance reaches the pair caps. The benchmark's 60 comparison rows are
not included in this combined CSV.}
\label{tab:corpus}
\end{table}

The HC3 JSONLs retain questions and answer lists but no acquisition log,
upstream revision, or complete reuse documentation. Those records need to be
restored before a fully documented corpus release. Questions establish the
pairing structure; the detector receives answer text only. Neither the
three-domain HC3 subset nor the news/magazine benchmark covers the variety
of human expertise, language proficiency, editing, and authorship practices
needed for broad deployment claims.

\subsection{Leakage control and evaluation splits}
The intended unit of separation is a connected cluster of related documents,
including shared source material, prompt variants, authors, and duplicates.
The actual implementations provide narrower protection. The original pilot
groups by article-topic pair, not publication. V3 separates human source
articles but uses different processing for human and generated prose.
The modern benchmark keeps numeric prompt groups, and therefore the audited
article pairs, within a single split. The V11 builder namespaces HC3 question
groups and reports zero shared prompt IDs between training and validation.
It prints that count rather than establishing every possible leakage relation.

\textbf{The revision audit finds exact-text leakage.} An identical human
medical answer appears under HC3 groups \code{medicine:416} and
\code{medicine:695} in training and \code{medicine:1124} in validation.
The text fields are byte-identical, before normalization. The duplication
already occurs in paired V10 and survives V11 construction. Zero overlapping
group IDs therefore does not justify the description ``leakage-free.''
The result is a concrete failure of grouping by question alone to capture
repeated answers.

This finding does not establish how much the duplicate affected the saved
metrics. Removing the validation row would not undo its relationship to
training, repeated development, or threshold selection; reporting modified
counts under the original experiment's name would be misleading. We preserve
the original results and require a newly versioned, deduplicated development
cycle with connected-component splitting. Near duplicates and repeated
source-specific language may create further dependencies not measured by
this exact-match audit.

Statistical separation also differs from informational separation. The same
60-row benchmark comparison partition was examined throughout model
development. Although it remains named \code{test} in several files, it is
not now an untouched final evaluation. V11 also incorporates benchmark
training and validation material, making that benchmark a development source
rather than independent external validation. The previously inspected pilot
test articles likewise became development material in the second cycle.
No fresh custodian-held final set is supplied for V11.

\subsection{Estimands and uncertainty}
A future prevalence study would estimate $p_A(t)$, $p_Q(t)$, and $p_S(t)$
separately. With audited labels and sampling weights $w_i$, a design-based
slop estimate would take the form
\begin{equation}
\widehat p_S(t)=\frac{\sum_{i\in t}w_iS_i}{\sum_{i\in t}w_i}.
\end{equation}
For a fixed screen with positive fraction $q$, sensitivity $r$, false-positive
rate $f$, and true prevalence $p$,
\begin{equation}
q=rp+f(1-p),\qquad \widehat p=\frac{q-f}{r-f}.
\label{eq:correction}
\end{equation}
The identity follows by conditioning on the true class. Its application
requires error rates that transfer to the sampled population. For $S$, those
are the \emph{joint screen's} errors, not an authorship model's errors. When
$r$ is close to $f$, the correction is unstable. Values outside $[0,1]$ should
prompt investigation rather than silent clipping.

No value of $p$ is estimated here. Class balance is a corpus-construction
choice, not a population prevalence, and V11 has no quality head. The original
proposal's 2,000 publisher-cluster bootstrap replicates and annual contrast
$p_S(2025)-p_S(2022)$ remain prospective. Any renewed longitudinal study must
fix its archive windows and sampling rules, account for extraction losses
and composition changes, and propagate uncertainty in period-specific error
rates. A positive growth estimate would not establish a majority; a
page-majority estimate would not establish exposure dominance.

\section{Detection Model and How the Code Works}
\subsection{Transparent features}
The original \code{legacy/detector.py} extracts ten interpretable surface
features: log token count, unique-token fraction, mean token length, mean
sentence length, sentence-length coefficient of variation, repeated-trigram
fraction, repeated-sentence fraction, digit-token fraction, URL count per
token, and selected punctuation count per token. The approximations are
deliberately simple. For example, splitting sentences at punctuation can
misread abbreviations, while repeated technical terms may be informative
rather than redundant. These features do not check facts or reveal provenance.

V11 replaces that representation with two learned TF--IDF vocabularies.
It contains no handcrafted quality, semantic-truth, perplexity, or
language-model-internal features. Table~\ref{tab:v11config} records the
configuration in the trainer, result JSON, and inspected serialized metadata.
The word and character vocabularies and inverse document frequencies are
fitted on training texts only; validation texts are transformed with those
fitted objects.

\begin{table}[htbp]
\centering\small
\begin{tabularx}{\linewidth}{lXX}
\toprule
Setting & Word block & Character block \\
\midrule
Analyzer & \code{word} & \code{char_wb} \\
$n$-grams & 1--3 words & 2--6 characters \\
Maximum features & 50,000 & 80,000 \\
Minimum document frequency & 2 training documents & 2 training documents \\
Term frequency & Sublinear & Sublinear \\
Vector normalization & L2 within block & L2 within block \\
Case handling & Lowercase & Lowercase \\
\midrule
Classifier & \multicolumn{2}{l}{\code{LinearSVC}, $C=3.0$, no class weighting} \\
Decision threshold & \multicolumn{2}{l}{$\tau=-0.364734947384451$} \\
\bottomrule
\end{tabularx}
\caption{V11 settings. The saved classifier expects 130,000 input coordinates.
Separate block normalization is followed by sparse concatenation, not an
additional normalization of the concatenated vector.}
\label{tab:v11config}
\end{table}

Before vectorization, the implemented normalization replaces HTML-like
substrings matched by \code{<[^>]+>} with spaces, removes double asterisks,
and removes isolated single asterisks matched by
\verb|(?<!\w)\*(?!\w)|. Curly double quotes become straight double quotes;
curly single quotes/apostrophes become straight apostrophes; en and em dashes
become hyphens. Repeated whitespace is collapsed and the result is stripped.
The vectorizers then lowercase text. This is a limited regex transformation,
not a complete HTML or Markdown parser: underscore emphasis and some
single-asterisk constructions are not removed. It does not remove all
formatting differences or normalize their semantic consequences.

For a retained feature $j$ occurring $c_{ij}>0$ times in document $i$, the
default smoothed, sublinear weighting is
\begin{equation}
v_{ij}=(1+\log c_{ij})\left(1+\log\frac{N+1}{d_j+1}\right),
\label{eq:tfidf}
\end{equation}
where $N$ and $d_j$ refer to the training corpus and its document frequency.
Absent features receive zero. Each block is L2-normalized. The word analyzer
uses the default token pattern retaining tokens of at least two word
characters; \code{char_wb} constructs character sequences within whitespace
word boundaries with boundary padding. Thus both lexical content and local
orthographic/formatting patterns can affect the score. These are potentially
useful predictors and potential shortcuts, not direct measurements of origin.

\subsection{Separate origin and quality classifiers}
The distinction in this heading belongs to the original design and remains
scientifically necessary, even though V11 trains only the origin classifier.
For the historical ten-feature baseline, training-only standardization is
\begin{equation}
z_{ij}=(x_{ij}-\mu_j)/s_j,
\end{equation}
with $s_j=1$ for effectively constant coordinates. The implemented logistic
heads have the form
\begin{equation}
h_k(\mathbf z)=\sigma(b_k+\mathbf w_k^\top\mathbf z),
\quad k\in\{A,Q\},\qquad \sigma(u)=\frac{1}{1+e^{-u}}.
\end{equation}
The original objective, with an unpenalized intercept, is
\begin{equation}
\mathcal L_k=\frac1n\sum_{i=1}^{n}
\left[\log(1+e^{\eta_{ik}})-y_{ik}\eta_{ik}\right]
+\frac{0.1}{2}\lVert\mathbf w_k\rVert_2^2,
\qquad \eta_{ik}=b_k+\mathbf w_k^\top\mathbf z_i.
\end{equation}
The code uses a stable logistic loss and an analytic gradient with SciPy's
L-BFGS-B optimizer. Pilot tuning changes the penalty and feature subset.
No independently trained quality head follows from implementing this interface.

V11 instead fits a supervised linear support vector classifier to the
concatenated TF--IDF vector $\mathbf u_i$. Its score is
\begin{equation}
f(\mathbf u_i)=b+\boldsymbol\beta^\top\mathbf u_i.
\label{eq:svm}
\end{equation}
Serialized metadata records an L2 penalty, squared-hinge loss, an intercept,
no class weighting, and $C=3.0$. The training code does not set a random seed
for the estimator and otherwise relies on its library defaults. The saved
metadata identifies scikit-learn 1.9.0. These details constrain reproducibility;
they are not a reason to infer that another software version will produce
byte-identical coefficients.

Both model families are small supervised classifiers. Fitting their weights
and selecting hyperparameters is not fine-tuning a language model. V11 neither
queries a generator nor reads hidden token probabilities, metadata from a
webpage, or Internet URLs. It evaluates only supplied text.

\subsection{Thresholds and explanations}
For V11, \code{train_combined_v11.py} searches the unique validation scores.
It maximizes validation recall subject to empirical validation FPR at most
5\%, then prefers lower FPR and finally a higher threshold. The fitted model
and vectorizers are fixed during this selection. The resulting rule is
\begin{equation}
\widehat A_i=\mathbb{1}\{f(\mathbf u_i)\geq\tau\},
\qquad m_i=f(\mathbf u_i)-\tau.
\label{eq:decision}
\end{equation}
The margin $m_i$ measures signed distance in score units from the chosen
operating point. It is not a probability. A score of 0.80 does not mean an
80\% probability of AI authorship, and a negative score can still produce
the output \code{AI} when the threshold is negative.

This empirical constraint is not a deployment guarantee. Its calculation
uses the same validation labels later summarized as validation performance.
The threshold is therefore selected rather than independently tested there.
Neither the model nor its threshold has been validated under every source,
topic, length, or generator shift. The number of validation negatives, their
dependencies, and the observed cross-split duplicate further limit inference.

The historical two-head design would combine decisions as
\begin{equation}
\widehat S_i=\mathbb{1}\{h_A(\mathbf z_i)\geq\tau_A\}
\mathbb{1}\{h_Q(\mathbf z_i)\geq\tau_Q\}.
\end{equation}
Even those logistic scores are not demonstrated to be calibrated, and their
product is not a validated slop probability. The original implementation
reports large feature contributions to each logit and returns a null joint
review flag when quality is untrained. Such contributions explain arithmetic,
not authorship. V11's CLI does not implement that explanation interface or
the joint decision; it reports a score, threshold, margin, word count, and
binary class string.

\subsection{Implementation boundary}
The original baseline abstains below 80 tokens. V11 does not inherit that
behavior: its CLI rejects empty files but imposes no validated minimum-length
rule and performs no out-of-distribution abstention. The HC3 builder's
80-word eligibility filter is a data-construction rule, not an inference
safeguard. A short text can therefore receive a V11 class string without
evidence that the operating point is appropriate for its length.

The current model bundle includes fitted vectorizers, classifier, threshold,
and experiment metadata. The compressed portable bundle retains the
vectorizers, classifier, and threshold; the release script supplies the same
normalization function locally. Neither requires the training corpus for
ordinary inference, but reconstructing training requires that corpus and its
software environment. The original bundle also references a normalization
function by its Python serialization name, which is an additional portability
consideration. Model serialization should be loaded only from trusted sources.

Path changes in the reorganized repository affect execution. The root
\code{detector_v11.py} expects a model in the working directory, whereas
the uploaded model is under \code{models/current}. Several builders expect
raw files at old root-level locations and trainers write outputs there.
The release command in the appendix uses the layout actually supplied.
These path issues do not invalidate recorded arithmetic, but they prevent
describing every README command as a verified turnkey reproduction.

\section{Results and Evaluation Plan}
\subsection{External findings versus original results}
Table~\ref{tab:status} separates the kinds of evidence. Reported model
metrics come from their archived JSONs. New audit computations are limited
to corpus reconstruction, identity checks, descriptive counts, and metric
arithmetic unless explicitly stated otherwise. No new model was trained and
no final evaluation was opened during manuscript revision.

\begin{table}[htbp]
\centering\small
\begin{tabularx}{\linewidth}{l>{\raggedright\arraybackslash}X>{\raggedright\arraybackslash}X}
\toprule
Evidence & Available record & What it does not establish \\
\midrule
Published prevalence studies & External literature & This project's slop prevalence \\
Software verification & Historical smoke-check records & Real-world classifier performance \\
Pilot and diagnostic cycles & Texts, code, weights, results & Modern-web generalization \\
Modern benchmark experiments & Grouped data, search and result JSONs & Untouched final evaluation after reuse \\
V11 validation & Combined CSV and aggregate metrics & Independent threshold assessment \\
External sanity checks & README observations; one supplied text & A representative external benchmark \\
Revision audit & Reconstructed membership, hashes, arithmetic & Provenance certification or inference replay \\
Quality and prevalence study & Rubric and prospective protocol & Completed annotations or population estimates \\
\bottomrule
\end{tabularx}
\caption{Evidence status of the current repository. The word ``test'' in a
filename does not establish independence from development.}
\label{tab:status}
\end{table}

\subsection{Observed software-verification results}
The historical record reports 21 initial smoke assertions, 36 additional
pilot checks, and 38 second-cycle statistical/workflow checks. These used
in-session or explicitly synthetic fixtures and are not a bundled independent
detector benchmark. Their successful completion should not be expressed as
classification accuracy on human and generated writing. The archived
verification records also describe earlier LaTeX builds, not the current
manuscript's scientific validity.

For this revision, \code{docs/audit_repository.py} fingerprints 171 research
inputs, inspects every processed CSV, compares the raw benchmark's texts and
labels with its processed version, reconstructs HC3 pairs from all three raw
files, and reconstructs combined V11 membership. Both HC3 reconstructions
match. It also checks 133 metric blocks against their confusion counts with
no arithmetic discrepancy. This last check cannot validate AUROC, which
requires score ordering, or establish that the underlying predictions are
correct. It does establish that the stored thresholded metrics agree with
the stored counts.

The revision environment lacks scikit-learn and joblib; no packages were
installed. Saved-model inference, AUROC, and row-level subgroup results were
therefore not independently reproduced. Non-executing serialized-metadata
inspection corroborates V11's vectorizer and classifier settings, but is
not a substitute for scoring the corpus. Recorded results are identified
as such throughout.

\subsection{Original pilot: ranking without detection}
The 1880 pilot allocates six topic pairs to training, three to validation,
and three to its original test. Eight configurations compare all ten
features with a subset omitting log token count and mean sentence length,
at L2 penalties 0.01, 0.1, 1, and 10. All eight achieve validation AUROC 1.0,
recall 1.0, and FPR 0 on six documents. The recorded tie-break selects the
eight-feature configuration with penalty 10, not a statistically demonstrated
winner. Its original test has three true negatives, three false negatives,
and no positive predictions. The baseline gives the same decisions.

Both fitted models nevertheless have test AUROC 1.0. For the selected model,
historical scores range from approximately 0.2776 to 0.4731 and generated
scores from 0.5146 to 0.5372; all are below the validation-selected threshold
of approximately 0.5400. The model preserves the ranking while missing every
generated test article. Its accuracy is 0.5, recall and F1 are zero, and
precision is undefined. An always-negative control has the same decisions
but AUROC 0.5. The experiment therefore distinguishes a useful ranking from
a useful operating point.

The observed FPR of zero is weak reassurance with three human negatives.
Even under independent Bernoulli sampling, its one-sided 95\% upper bound
would be $1-0.05^{1/3}\approx0.632$. This is an illustrative small-sample
calculation, not a credible population guarantee for a shared-source corpus.
The original threshold and failed decisions remain part of the record after
those articles are reassigned to development.

\subsection{Model evolution from V2 through V11}
V2 moves from handcrafted features to word 1--3-gram and character 3--5-gram
TF--IDF with logistic regression. Its 24-document pilot record reports
perfect thresholded performance in the 12/6/6 partitions. V3 retains that
family of representation on the separate 60-document historical expansion
and also reports perfect results. These are not independent confirmations
of the original pilot: its data had already been examined, and the expansion
remains historically confounded.

V3 adds a particularly visible processing asymmetry. The human chunk builder
tokenizes prose and rejoins tokens with spaces, removing punctuation, whereas
the fixed generated articles are retained as ordinary prose. Origin is thus
correlated with preprocessing. High performance can reflect that difference
as well as era, source, or writing style; the experiment does not isolate
their contributions.

The modern benchmark then becomes the principal development resource.
The archived benchmark V2 logistic model reaches comparison accuracy 0.9333,
with three false positives and one false negative. A 21-configuration,
three-fold GroupKFold search compares word, character, and combined features
over seven regularization values. The selected combined logistic model
uses $C=10$ and records comparison accuracy 0.9667. The search metadata's
older textual group description does not override the actual grouped CSV
used by the code.

V4 broadens the search to 48 settings, comparing logistic regression and
linear SVMs with word, character, combined, and stylometric-augmented
representations. The selected model uses character \code{char_wb} 3--5-grams,
40,000 features, and $C=10$. V5 evaluates 135 settings spanning character
ranges, combined features, regularization, model family, and class weights.
Its selected character 2--6-gram SVM uses 60,000 features and $C=10$ without
class weighting. V5's grouped search refits features inside each fitting
fold, but chooses a threshold and reports decision metrics on that same
fold's validation labels. Those operating-point metrics are selection
diagnostics, not nested estimates of a frozen decision procedure.

\begin{table}[htbp]
\centering\small
\begin{tabular}{lrrrrrr}
\toprule
Model / assessment & $n$ & Accuracy & Recall & FPR & FNR & AUROC \\
\midrule
Pilot selected / original test & 6 & .5000 & .0000 & .0000 & 1.0000 & 1.0000 \\
V2 / pilot comparison & 6 & 1.0000 & 1.0000 & .0000 & .0000 & 1.0000 \\
V3 / 1900 comparison & 16 & 1.0000 & 1.0000 & .0000 & .0000 & 1.0000 \\
Benchmark V2 / comparison & 60 & .9333 & .9667 & .1000 & .0333 & .9967 \\
Optimized logistic / comparison & 60 & .9667 & 1.0000 & .0667 & .0000 & .9989 \\
V4 / comparison & 60 & .9667 & 1.0000 & .0667 & .0000 & 1.0000 \\
V5 / comparison & 60 & .9833 & 1.0000 & .0333 & .0000 & .9989 \\
V6 / full-text comparison & 60 & .8500 & 1.0000 & .3000 & .0000 & 1.0000 \\
V7 / comparison & 60 & .9833 & 1.0000 & .0333 & .0000 & .9989 \\
V8 / comparison & 60 & .9833 & 1.0000 & .0333 & .0000 & .9989 \\
V11 / combined validation & 620 & .9742 & .9935 & .0452 & .0065 & .9991 \\
\bottomrule
\end{tabular}
\caption{Archived experiment results, rounded to four decimals; FNR is
derived as $1-\mathrm{recall}$. Different corpora and evaluation roles must
not be ranked as a single leaderboard. The benchmark comparison rows are
reused across development, and V11's row is validation, not a new test.}
\label{tab:evolution}
\end{table}

V6 addresses document length by forming non-overlapping 180--260-word chunks
from the benchmark's training documents. It obtains 474 training chunks
(241 human, 233 generated) and 160 chunks from the separate validation
documents (82 human, 78 generated). Chunk creation respects the existing
document split, but multiple chunks remain dependent observations. The
short-validation threshold gives three human false positives and one AI
false negative. On the original full-document comparison set it gives nine
false positives among 30 human texts and no false negatives among 30 generated
texts. Thus AUROC 1.0 coexists with FPR 0.30. Improved sensitivity on the
reported short external examples comes at a substantial cost under this
different length distribution.

V7 combines V5 and V6 margins, using the first 260 words for the short model
on longer documents. Twenty-one mixture weights are compared on validation.
The selected weights are 1.0 for V5 and 0.0 for V6, with a threshold of zero
on the resulting V5 margin. The selected ``ensemble'' is therefore V5's
decision rule, not evidence that combining the two models improved transfer.

V8 applies the normalization described above to character features. It
records one human false positive and no AI false negatives on the 60-row
comparison set, the same confusion counts as V5, although a matching count
does not prove that every prediction is identical. The README reports that
normalization did not resolve external generalization. Its outcome motivates
broader data coverage rather than a claim that formatting is irrelevant.

A V9 preparation step selects 227 benchmark training/validation documents
of 350--1,000 words; no fitted V9 result is supplied. The next material
change is HC3 acquisition, followed by V10 pairing and V11 combination.
No separate fitted V10 result is supplied either. V11 changes training
composition, feature combination, feature limits, document-frequency
filtering, and regularization relative to V8. Its reported external
improvement is consistent with the value of broader data, but this is not
a controlled ablation establishing diversity as the sole cause.

\subsection{V11 validation results and unresolved subgroup attribution}
Table~\ref{tab:v11metrics} preserves the aggregate values in
\code{models/current/hc3_v11_combined_results.json}. The training record
has no thresholded errors in 2,248 texts and AUROC 1.0; these are fit-set
statistics, not generalization evidence. In validation, 14 of 310
human-labeled texts are flagged and two of 310 generated texts are missed.
The recorded threshold satisfies the empirical 5\% constraint on that
selection set. It does not constrain each source separately.

\begin{table}[htbp]
\centering\small
\begin{tabular}{lr}
\toprule
Quantity & Recorded value \\
\midrule
Validation texts & 620 \\
Accuracy & 0.9741935483870968 \\
Precision & 0.9565217391304348 \\
Recall & 0.9935483870967742 \\
F1 & 0.9746835443037974 \\
False-positive rate & 0.04516129032258064 \\
AUROC & 0.9991155046826222 \\
True negatives / false positives & 296 / 14 \\
False negatives / true positives & 2 / 308 \\
\bottomrule
\end{tabular}
\caption{Exact aggregate V11 validation values from the archived JSON.
The derived FNR is $2/310\approx0.006452$. Numerical precision records the
artifact, not the certainty of population performance.}
\label{tab:v11metrics}
\end{table}

Source-specific error rates would be more informative than an aggregate
alone, but the upload does not retain the necessary output. The V11 trainer
prints metrics by origin and source; its saved result JSON contains only
aggregate training and validation metrics. There is no associated stdout
ledger or per-document validation score file. Table~\ref{tab:subgroups}
therefore reports auditable denominators without inventing source errors.

\begin{table}[htbp]
\centering\small
\begin{tabular}{lrrl}
\toprule
Validation component & Human & AI & Source FPR / recall \\
\midrule
HC3 finance & 150 & 150 & Not retained in V11 JSON \\
HC3 medicine & 126 & 126 & Not retained in V11 JSON \\
HC3 open QA & 4 & 4 & Not retained in V11 JSON \\
Modern benchmark & 30 & 30 & Not retained in V11 JSON \\
\bottomrule
\end{tabular}
\caption{V11 subgroup denominators. Missing subgroup errors are not zeros.
The small open-QA subgroup cannot support a stable source-level inference.}
\label{tab:subgroups}
\end{table}

The project handoff also reports approximate HC3 figures: finance human FPR
0.057 and AI recall 0.9962; medicine FPR 0.0159 and recall 1.0; open-QA FPR
0.25 with four human examples. The finance values cannot be produced by
integer error counts with 150 observations per label. They are arithmetically
compatible with $30/526$ and $524/526$ in the \emph{uncapped paired}
validation corpus, which contains 526 finance examples per label. That
compatibility does not identify the model or run that generated them.
Medicine and open-QA denominators are unchanged by capping and do not resolve
the attribution. We preserve this discrepancy as an unresolved record,
not a V11 source-performance table.

The handoff's reported contrast between benchmark and HC3 validation likewise
needs its original output or a saved-model replay. It must not be suppressed
by implying uniform performance, but it cannot be quantified from aggregate
counts alone. A reproducible breakdown must use the archived model,
normalization, frozen threshold, exact combined CSV, and explicit subgroup
membership. No threshold should be changed during that reconstruction.

\subsection{External sanity checks and formatting dependence}
The README records the four V11 outputs in Table~\ref{tab:external}. It also
reports that V5 missed \code{article.txt}, V6 improved sensitivity to short
external examples, and V8 did not resolve transfer. These observations
motivated development and should remain visible. They are weaker evidence
than an archived, independently labeled evaluation with preserved scores.

\begin{table}[htbp]
\centering\small
\begin{tabularx}{\linewidth}{llX}
\toprule
Example & Reported V11 output & Evidence available in this upload \\
\midrule
\code{article_human.txt} & HUMAN & README observation; text absent \\
\code{article_ai.txt} & AI & README observation; text absent \\
\code{article_mixed.txt} & AI & README observation; text absent; mixed origin \\
\code{article.txt} & AI & Text present under \code{samples}; 175 words \\
\bottomrule
\end{tabularx}
\caption{Project-level development sanity checks, not a four-document
benchmark independently reproduced in this revision. Source/generation
provenance and per-file score logs are incomplete.}
\label{tab:external}
\end{table}

An AI output for a mixed-authorship text cannot be counted as a correct
binary-origin decision without first defining the target. The existing
\code{article.txt} concerns automated web traffic, but its claims are sample
input, not evidence supporting this paper's prevalence argument. Its
175-word length is substantially below the benchmark means. The README's
approximate 150--300-word description of external examples cannot be
verified for the three missing files. None of these observations estimates
an external accuracy rate, and no 100\% accuracy claim follows from them.

The README describes a feature-weight audit that found prominent HTML,
Markdown, and quotation patterns. A dedicated weight-audit output is not
retained. The revision can nevertheless verify label-correlated formatting
directly in the stored benchmark. Table~\ref{tab:formatting} gives document
counts for specified patterns across all 300 rows. These are retrospective
descriptions of an already inspected corpus, not additional validation.

\begin{table}[htbp]
\centering\small
\begin{tabular}{lrr}
\toprule
Pattern present in a document & Human ($n=150$) & AI ($n=150$) \\
\midrule
Any HTML-like tag & 83 & 72 \\
Double asterisks & 0 & 23 \\
\code{<br>} variants & 30 & 60 \\
Curly double quotation marks & 144 & 65 \\
Straight double quotation marks & 11 & 86 \\
\bottomrule
\end{tabular}
\caption{Audit counts from the raw text fields in the grouped benchmark.
Categories overlap; they are not exhaustive or mutually exclusive. Exact
regex definitions are retained in the audit script.}
\label{tab:formatting}
\end{table}

The relevant mechanism is straightforward: a classifier can use a pattern
that predicts the dataset label without that pattern being characteristic
of authorship in another population. The counts show opportunities for
such behavior, not that any one feature caused all errors. Normalization
addresses some visible differences but can leave topic vocabulary, source
conventions, editorial style, and other artifacts intact. V8's reported
transfer failure is consistent with that broader problem; it does not
identify its unique cause.

\subsection{What the completed experiment must report}
A confirmatory successor must evaluate a model and threshold frozen before
access to new final labels. It should report TP, FP, TN, FN, FPR, FNR,
precision, recall, F1, accuracy, AUROC, and any abstention coverage, together
with denominators and uncertainty. The current V11 point estimates should
not be supplemented with invented confidence intervals. A credible interval
analysis must account for question/article clustering, threshold selection,
model selection, and the actual sampling frame; an ordinary independent-row
binomial interval would miss important dependencies here.

Threshold stability needs repeated development partitions with a separate
calibration stage, followed by topic-, source-, time-, and generator-disjoint
assessment. Length strata, formatting interventions, paraphrasing, human-edited
AI text, and documented proficiency groups should be included where provenance
and consent permit. A shift-specific threshold can be studied prospectively,
but not introduced after inspecting final errors to obtain a preferred result.
RAID remains a possible origin-detection challenge resource, not a
quality-labeled slop benchmark \cite{raid}.

Balanced evaluation also does not determine deployment precision. For a
purely hypothetical prevalence of 5\%, sensitivity of 80\%, and FPR of 5\%,
1,000 documents yield expected counts of 40 true positives and 47.5 false
positives. Precision would be about 45.7\%. This is arithmetic illustrating
the base-rate issue, not an observed result for V11. Population measurement
requires the additional assumptions in Equation~\ref{eq:correction}.

\section{Second Development Cycle: Diagnostics and Readiness}
\subsection{Scope and separation of evidence}
The second cycle preceded the modern benchmark and V11 work. It reused all
24 original pilot articles as development material to investigate the failed
operating point. Its archival readiness files correctly describe what had
not been obtained \emph{at that stage}; statements that contemporary data
were absent are no longer a description of the entire current repository.
They also do not imply that the newer corpora meet the earlier quality,
rights-review, or final-custody requirements.

The main diagnostic implementation is preserved in
\code{legacy/cycle2_diagnostics.py}; \code{docs/research_cycle.py} supplies
the supplementary repeated-split analysis. The latter still imports the
old detector name and uses old root-relative paths. It is neither V11's
trainer nor a current, completed confirmatory evaluation. Its retained
outputs, not an assumed successful rerun in the reorganized checkout, are
the source of the diagnostic results summarized here.

\subsection{Exploratory failure analysis and controls}
The primary analysis uses 12 leave-one-topic-pair-out assessments. In each
fold, one pair is assessed, three select a threshold, and eight fit the
model. Training-only normalization is recomputed within each fold. The
original ten-feature setting gives FNR 0.25 and the selected eight-feature
setting gives FNR $1/3$; both have FPR zero and mean within-fold AUROC 1.0.
A length-only control has FPR $1/6$ and FNR $2/3$, while removing repetition
features leaves the original setting's aggregate decisions unchanged.
Scores from different fits are not pooled into a single AUROC.

The original model's thresholds range from approximately 0.6420 to 0.9373,
with median 0.8015. The selected setting ranges from 0.5089 to 0.5310,
with median 0.5165. Its smaller numerical movement is not proof of greater
stability: different penalties alter score scales, and the setting still
misses more generated articles in these folds. Conditional topic-pair
bootstrap intervals for FNR are $[0,0.50]$ and approximately
$[0.083,0.583]$, respectively. These archived intervals exclude refitting
uncertainty and do not estimate modern-web error.

The supplementary analysis evaluates 100 fixed-seed 6/3/3 topic-pair
partitions, five feature/penalty settings, and three threshold policies.
There are 500 fits for those comparisons and 100 additional permutation
fits, not 600 independent experiments. For the selected pilot setting,
the original threshold rule has mean FNR 0.240 and mean FPR 0.000. Using
the next representable value above the largest calibration-negative score
reduces mean FNR to zero but raises mean FPR to 0.260. A separating midpoint,
with fallback to the original rule, gives mean FNR about 0.007 and FPR 0.040.
These counterfactual policies describe an exploratory trade-off; they do not
replace the original failed test result.

Formatting probes on the fixed pilot models also matter. Uppercasing and
collapsing whitespace leave scores unchanged, while adding a hyphen before
each sentence changes five of 24 baseline decisions and four of 24 selected
model decisions. Maximum absolute score changes are about 0.1541 and 0.01849.
A digit-to-``number'' substitution changes content and must not be called
a meaning-preserving formatting test.

The 99-permutation leave-pair-out analysis gives a Monte Carlo value of
0.01 for mean within-fold AUROC. The separate 100-permutation analysis on
the original assessment split gives a tail fraction about 0.099. They use
different statistics and allocations; reporting only the more favorable
value would misrepresent the record. Both are exploratory association checks
under exchangeability assumptions, not evidence that origin rather than
era, source, or register caused the association.

The immediate failure mechanism is identifiable: generated assessment
scores remain below a conservative validation-derived cutoff despite
preserved rank ordering. Its underlying cause is not isolated. Historical
era, publication, and generating assistant are aligned with origin, so a
source-disjoint comparison is not identifiable within this tiny corpus.
Repeated splitting cannot create missing source diversity.

\subsection{Uncertainty and prospective methodology}
The archived prospective protocol fixes the original ten-feature logistic
model, not V11, and proposes negative-only Neyman--Pearson order-statistic
calibration. Such methods separate calibration from fitting under explicit
assumptions \cite{tong2018}. The retained implementation uses per-head
target FPR 0.05 and violation probability 0.025, with a planned minimum of
120 independent negative clusters per head. An order-statistic bound requires
independent, exchangeable calibration negatives and a scoring model fixed
before calibration. It is not a guarantee under distribution shift and is
not the empirical threshold search actually used for V11.

That distinction prevents a tempting but invalid inference: the existence
of a conservative calibration function elsewhere in the repository does
not confer its guarantees on the current SVM. The local method snapshot
is also not an externally timestamped preregistration. It neither freezes
V11 retrospectively nor converts its validation results into confirmatory
evidence. A new V11-based protocol would require a versioned amendment
before any new final evaluation.

\subsection{Annotation, custody, and release status}
The archived rubric proposes at least 100 independently double-rated
documents, raw binary agreement at least 0.80, a lower 95\% cluster-bootstrap
kappa bound at least 0.60, and at least 95\% resolved quality decisions.
These are design choices, not established universal acceptance standards.
The actual annotation audit has zero eligible documents and null agreement
statistics. New benchmark origin guesses do not satisfy these requirements.
There is no completed quality agreement analysis to report.

The prospective workflow records human review, rights and provenance checks,
model/input hashes, and an independent custodian's commitment for final
data. Software can check the supplied records but cannot establish that
raters were truly independent or that final data remained unseen. No
completed V11 evaluation satisfying that workflow is supplied. Original
pilot and later benchmark materials remain development evidence, and a
fresh final corpus is still required.

\section{Limitations, Ethics, and Next Steps}
\textbf{Coverage and confounding.} V11 covers a limited mixture of finance,
medicine, open-domain answers, and news/magazine writing. The benchmark
conditions do not represent every generator, decoding procedure, future
model, or degree of human editing. Historical experiments are further
confounded by era, source, and processing. Neither training size nor
pair-level balance resolves those limitations. Authorship can correlate
with topic, proficiency, translation patterns, formatting, and source
conventions without the classifier learning a portable origin signal.

\textbf{Evaluation dependence.} Development repeatedly inspected the
benchmark comparison set and external examples. V11's aggregate comes
from the threshold-selection set and includes an exact cross-split text
duplicate. Some source strata are extremely small: four human open-QA
examples allow an empirical FPR change of 0.25 from a single error. The
source-specific errors are not retained for the combined model, and the
handoff numbers cannot safely be attributed to that validation set. The
available record supports neither a stable source-wise FPR claim nor an
unbiased universal accuracy estimate.

\textbf{What the score cannot say.} A human document can be flagged and a
generated one can be missed. Heavy paraphrasing, translation, partial
assistance, and later editing may change both the appropriate target and
the feature distribution. V11 has no calibrated probabilities, no quality
head, and no validated short-text or out-of-distribution abstention rule.
Its lexical weights cannot establish authorship from text alone. No
population prevalence or annual trend follows from a balanced convenience
corpus and its validation scores.

\textbf{Reproducibility and provenance.} The audit verifies important
transformations but does not reproduce fitted-model inference in the
available environment. Missing external texts, row-level score logs,
acquisition/generation records, and reuse documentation limit independent
replication. Some script paths no longer match the repository layout.
These are research-package limitations, not details to conceal in a
favorable performance summary. The archived artifacts should remain intact
while a corrected, separately versioned pipeline is developed.

\textbf{Responsible use.} Outputs are screening and research signals for
manual review, not proof of cheating, plagiarism, origin, or misconduct.
They should not be used alone for academic misconduct accusations, punitive
moderation, employment decisions, publication rejection, legal conclusions,
or automatic content removal. The CLI's \code{AI}/\code{HUMAN} strings are
operational class names, not certified findings about an individual. False
accusations cannot be prevented merely by choosing a low empirical validation
FPR, and a low score does not certify human authorship.

\textbf{Collection and disclosure.} A future release should use permitted
public material or consented writing, record reuse permissions, minimize
personal information, and prefer aggregate findings over unnecessary author
identifiers or annotator comments. A source link is not itself a reuse
permission, and a historical Gutenberg notice does not authorize release
of unrelated contemporary texts. No new article bodies are redistributed
by this manuscript. An AI assistant was used to prepare and revise study
materials and this paper, and authored the original pilot's generated
articles. The later fixed generated examples are retained as development
materials, not represented as independent human writing. Venue-specific
disclosure and accountable human review remain necessary before submission.

The next cycle should first repair provenance and grouping, restore missing
inputs and per-document result logs, and preserve the V11 snapshot. After
constructing a genuinely new, independently documented corpus, developers
should conduct source/topic/generator and length controls without accessing
the final set. If quality remains a target, independent rubric annotation
and an agreement audit must precede quality-model fitting. Freeze the
features, calibration rule, model, primary metrics, and subgroup plan before
custodian-held final and challenge evaluation. Only after transferable error
rates and a probability sample exist should error-adjusted prevalence be
attempted. None of those requirements can be completed by editing prose.

\section{Conclusion}
The project's clearest result is not its highest accuracy value. It is the
sequence of ways a favorable score can fail to answer the intended question.
The initial pilot preserves ranking but fails at its chosen threshold. V5
performs strongly on a repeatedly inspected benchmark while missing a
reported short external example. V6 shifts the operating point at the cost
of substantially more benchmark human false positives. V8 addresses visible
formatting artifacts without resolving reported transfer failures. V11
adds paired HC3 data and shows improved behavior in the available project
notes, alongside strong aggregate validation performance.

That progression supports further study of training composition and
distribution shift, not a causal claim that one change solved detection.
The audit strengthens the contribution by identifying exact-text leakage,
ambiguous subgroup attribution, and missing external evidence rather than
leaving those limitations implicit. The present work is an exploratory
detector-development case study with an explicit route to stronger
evaluation. It does not establish a universally accurate AI detector,
a validated low-quality-content detector, or an Internet-wide ``slop''
majority. Those claims require evidence that this repository does not yet
contain.

\appendix
\section{Reproducing the Implementation}
\subsection{Evidence map and audit}
The revision's supporting files are \code{docs/PROJECT_AUDIT.md},
\code{docs/audit_repository.py}, and \code{docs/audit_results.json}.
The pre-revision manuscript is retained as \code{docs/main_pre_v11.tex};
it is a historical snapshot, not an alternative current paper. Table
\ref{tab:evidencemap} identifies the principal sources for the numerical
and methodological claims. Full input hashes and constructor expressions
are included in the machine-readable audit.

\begin{table}[htbp]
\centering\small
\begin{tabularx}{\linewidth}{lX}
\toprule
Claim & Repository evidence \\
\midrule
Pilot selection and failed test & \code{legacy/pilot/artifacts/selection.json}; \code{test_results.json}; \code{test_predictions.json} in the same directory \\
Historical V2/V3 & \code{models/archive/v2_results.json}; \code{v3_results.json}; \code{dataset/v3_corpus.csv} \\
Benchmark evolution & \code{models/archive/hc3_*results.json}; optimization/search JSONs and corresponding \code{training/} scripts \\
Original/modern distinction & \code{legacy/pilot/protocol.json}; \code{data/historical_1900/}; \code{data/benchmark/human_detectors.json} \\
HC3 pairing and V11 composition & \code{training/build_hc3_paired_v10.py}; \code{training/build_combined_v11.py}; their output CSVs \\
V11 model and aggregate metrics & \code{training/train_combined_v11.py}; \code{models/current/hc3_v11_combined_results.json}; saved model metadata \\
Second-cycle diagnostics & \code{legacy/cycle2/artifacts/diagnostics.json}; \code{legacy/cycle2/results/diagnostics.json} \\
Incomplete annotation & \code{legacy/cycle2/artifacts/annotation_audit.json}; blank quality fields in processed data \\
External observations & \code{README.md}; \code{samples/article.txt}; no complete per-file score ledger \\
\bottomrule
\end{tabularx}
\caption{Evidence map. Patterned filenames denote a family of retained
artifacts, not a claim that every file used the same dataset.}
\label{tab:evidencemap}
\end{table}

From the repository root, the non-training audit runs with the supplied
Python environment and pandas:
\begin{lstlisting}
python docs/audit_repository.py
\end{lstlisting}
It overwrites only its revision audit JSON, not models or experiment results.
It checks exact text duplication and supplied group membership, not every
near-duplicate or shared-author relationship. It does not retrieve links,
verify claims within articles, determine reuse rights, or infer missing labels.

\subsection{Inference, training, and software versions}
The release layout provides an inference-only route when the recorded
dependencies are available:
\begin{lstlisting}
cd release/V11_RELEASE
python3 detector_v11.py ../../samples/article.txt
\end{lstlisting}
The command must be launched from that folder because the model path is
working-directory relative. It describes the supplied layout; inference
was not rerun in the revision environment. The current and release portable
model copies have matching SHA-256 hashes.

The uploaded dependency file lists NumPy 2.5.2, SciPy 1.18.0,
scikit-learn 1.9.0, joblib 1.5.3, and pandas 3.0.5; the README records
Python 3.13.14. In contrast, the audit ran on Python 3.13.15, NumPy 2.4.1,
SciPy 1.17.0, and pandas 2.3.3, without scikit-learn or joblib. These are
distinct environments. No cross-version inference equivalence or exact
training replay is claimed.

The V11 training sequence is implemented by
\code{training/build_hc3_paired_v10.py},
\code{training/build_combined_v11.py}, and
\code{training/train_combined_v11.py}. Before replay, restore or explicitly
adapt their expected input/output locations in a separate experiment
directory. Do not run them blindly over preserved artifacts. Reproduction
should log library versions, raw and processed hashes, pair membership,
feature settings, threshold, model fingerprint, all row scores, and
source-level metrics. A data-cleaned successor needs a new identifier and
must not replace the V11 results reported here.

Build the current self-contained paper from the repository root with:
\begin{lstlisting}
latexmk -cd -pdf -interaction=nonstopmode -halt-on-error docs/main.tex
\end{lstlisting}
The bibliography is embedded. Tables require no external pilot figures or
obsolete root-relative inputs. The root-level \code{main.pdf} belongs to
the earlier project state; the revised output is \code{docs/main.pdf}.

\begin{thebibliography}{99}
\raggedright
\bibitem{pew}
S. Bestvater, A. Smith, C. TerBush, C. Baronavski, and J. Chavda.
\emph{How Much of the Internet Is Written With AI?}
Pew Research Center, August 20, 2026.

\bibitem{pewmethods}
S. Bestvater, A. Smith, C. TerBush, C. Baronavski, and J. Chavda.
\emph{Methodology: How Much of the Internet Is Written With AI?}
Pew Research Center, August 20, 2026.

\bibitem{dolezal}
J. Dolezal, S. Alam, M. Graham, and M. Bohacek.
\emph{The Impact of AI-Generated Text on the Internet}.
arXiv:2604.26965, version 1, April 14, 2026. Preprint.

\bibitem{raid}
L. Dugan, A. Hwang, F. Trhl\'{i}k, A. Zhu, J. M. Ludan, H. Xu,
D. Ippolito, and C. Callison-Burch.
RAID: A Shared Benchmark for Robust Evaluation of Machine-Generated Text
Detectors. \emph{Proceedings of ACL}, 2024.
doi:10.18653/v1/2024.acl-long.674.

\bibitem{tufts}
B. Tufts, X. Zhao, and L. Li.
A Practical Examination of AI-Generated Text Detectors for Large Language
Models. \emph{Findings of NAACL}, pp. 4839--4856, 2025.
doi:10.18653/v1/2025.findings-naacl.271.

\bibitem{liang}
W. Liang, M. Yuksekgonul, Y. Mao, E. Wu, and J. Zou.
GPT detectors are biased against non-native English writers.
\emph{Patterns}, 4(7):100779, 2023.
doi:10.1016/j.patter.2023.100779.

\bibitem{tong2018}
X. Tong, Y. Feng, and J. J. Li.
Neyman-Pearson classification algorithms and NP receiver operating characteristics.
\emph{Science Advances}, 4(2):eaao1659, 2018.
doi:10.1126/sciadv.aao1659.

\bibitem{m4}
Y. Wang et al.
M4: Multi-generator, Multi-domain, and Multi-lingual Black-Box Machine-Generated
Text Detection. \emph{Proceedings of EACL}, pp. 1369--1407, 2024.
doi:10.18653/v1/2024.eacl-long.83.

\bibitem{coauthor}
M. Lee, P. Liang, and Q. Yang.
CoAuthor: Designing a Human-AI Collaborative Writing Dataset for Exploring
Language Model Capabilities. \emph{Proceedings of CHI}, 2022.
doi:10.1145/3491102.3502030.

\bibitem{summeval}
A. R. Fabbri, W. Kry\'{s}ci\'{n}ski, B. McCann, C. Xiong, R. Socher, and D. Radev.
SummEval: Re-evaluating Summarization Evaluation.
\emph{Transactions of the Association for Computational Linguistics}, 9:391--409,
2021. doi:10.1162/tacl\_a\_00373.
\end{thebibliography}
\end{document}
